\begin{center}
	{
		\Large
		\textbf{Physics 225 - General Relativity 1\\[10pt] Course Information}
	}    
\end{center}
\addcontentsline{toc}{section}{Course Information}

\subsubsubsection*{Course Description}
	\paragraph{Description: } 
	Manifolds, modern differential geometry and tensor analysis; basic principles of general relativity;
	Einstein field equations and their mathematical properties; exact solutions; linearized theory; variational
	principles and conservation laws; equations of motion; gravitational waves; experimental tests.

\subsubsubsection*{Instructor Details}
	\begin{itemize}[]
		\item[] \textbf{Instructor: } Dr. Michael Francis Ian Vega
		\item[] \textbf{Email: } {\texttt{\url{ivega.upd.edu.ph}}}
		\item[] \textbf{Lecture Hours: } WF, 10:00 - 11:30 AM
		% \item[] \textbf{Classroom Number: } R109
	\end{itemize}

\subsubsubsection*{Course Requirements}
	\paragraph{Problem Sets} 

	Problem sets will be assigned throughout the semester, ideally one per week. These will be
	submitted electronically using a Google Form that I will be sending out close to the deadline. From each
	problem set a random selection of items will be graded. It is your responsibility to ensure that your work
	is neat and legible. If we cannot understand your work then you will not get any credit. These problem
	sets will make up the 30\% of your final grade.

	\paragraph{Written Exam} 
	There will be two written examns in the course. Each exam will constitute 30\% of the final grade.
	These are in-class, open-notes exams. You will have to designate one notebook as your 225
	notebook at the start of the semester. Only this notebook may be used to help you during the exams.

	\paragraph{Oral Exam} This will be scheduled during or shortly after Finals Week. This will make up the remaining 10\%
	of your final grade.


